\documentclass[11pt,a4paper]{article}
\usepackage[utf8]{inputenc}
\usepackage[margin=1in]{geometry}
\usepackage{hyperref}
\usepackage{listings}
\usepackage{xcolor}
\usepackage{graphicx}
\usepackage{enumitem}
\usepackage{fancyhdr}
\usepackage{titlesec}

% Code listing style
\definecolor{codegreen}{rgb}{0,0.6,0}
\definecolor{codegray}{rgb}{0.5,0.5,0.5}
\definecolor{codepurple}{rgb}{0.58,0,0.82}
\definecolor{backcolour}{rgb}{0.95,0.95,0.92}

\lstdefinestyle{mystyle}{
    backgroundcolor=\color{backcolour},   
    commentstyle=\color{codegreen},
    keywordstyle=\color{magenta},
    numberstyle=\tiny\color{codegray},
    stringstyle=\color{codepurple},
    basicstyle=\ttfamily\footnotesize,
    breakatwhitespace=false,         
    breaklines=true,                 
    captionpos=b,                    
    keepspaces=true,                 
    numbers=left,                    
    numbersep=5pt,                  
    showspaces=false,                
    showstringspaces=false,
    showtabs=false,                  
    tabsize=2
}

\lstset{style=mystyle}

% Header and footer
\pagestyle{fancy}
\fancyhf{}
\rhead{Realtime Threads \& Chat Application}
\lhead{Technical Documentation}
\rfoot{Page \thepage}

\title{\textbf{Realtime Threads \& Chat Application}\\
\large Technical Documentation for Interview Discussion}
\author{Full-Stack Development Project}
\date{\today}

\begin{document}

\maketitle
\tableofcontents
\newpage

\section{Executive Summary}

This is a full-stack realtime web application that implements a social media platform with threaded discussions and direct messaging capabilities. The project demonstrates proficiency in modern web technologies, real-time communication, database design, authentication systems, and API development.

\subsection{Key Highlights}
\begin{itemize}[leftmargin=*]
    \item \textbf{Custom Authentication:} Implemented JWT-based authentication with bcrypt password hashing, eliminating third-party dependencies
    \item \textbf{Real-time Features:} WebSocket-based chat and notifications using Socket.IO
    \item \textbf{Database Migration:} Successfully migrated from PostgreSQL to SQLite for simplified deployment
    \item \textbf{Full-Stack TypeScript:} Type-safe development across frontend and backend
    \item \textbf{Modern UI/UX:} Responsive design with Next.js 16 and shadcn/ui components
\end{itemize}

\section{Technology Stack}

\subsection{Frontend}
\begin{itemize}[leftmargin=*]
    \item \textbf{Framework:} Next.js 16.0.5 (React 19.2.0)
    \item \textbf{Language:} TypeScript 5
    \item \textbf{UI Components:} shadcn/ui (Radix UI primitives)
    \item \textbf{Styling:} Tailwind CSS 3.4.1
    \item \textbf{State Management:} React Context API (custom AuthProvider)
    \item \textbf{HTTP Client:} Axios 1.7.9
    \item \textbf{Real-time:} Socket.IO Client 4.8.1
    \item \textbf{Notifications:} Sonner (toast notifications)
\end{itemize}

\subsection{Backend}
\begin{itemize}[leftmargin=*]
    \item \textbf{Runtime:} Node.js with Express 4.21.2
    \item \textbf{Language:} TypeScript 5
    \item \textbf{Database:} SQLite with better-sqlite3 11.0.0
    \item \textbf{Authentication:} jsonwebtoken 9.0.3 + bcrypt 5.1.1
    \item \textbf{File Upload:} Multer 1.4.5-lts.1 (local storage)
    \item \textbf{Real-time:} Socket.IO 4.8.1
    \item \textbf{Environment:} dotenv for configuration
\end{itemize}

\subsection{Development Tools}
\begin{itemize}[leftmargin=*]
    \item ts-node-dev for backend hot-reload
    \item ESLint for code quality
    \item Docker Compose support
\end{itemize}

\section{System Architecture}

\subsection{High-Level Architecture}
The application follows a client-server architecture with real-time capabilities:

\begin{verbatim}
┌─────────────────────────────────────────────────────────┐
│                    Client (Port 4000)                   │
│  ┌──────────────┐  ┌──────────────┐  ┌──────────────┐  │
│  │   Next.js    │  │  Socket.IO   │  │    Axios     │  │
│  │   Pages      │  │   Client     │  │   API Client │  │
│  └──────────────┘  └──────────────┘  └──────────────┘  │
└─────────────────────────────────────────────────────────┘
                            │
                    HTTP + WebSocket
                            │
┌─────────────────────────────────────────────────────────┐
│                   Server (Port 5000)                    │
│  ┌──────────────┐  ┌──────────────┐  ┌──────────────┐  │
│  │   Express    │  │  Socket.IO   │  │     JWT      │  │
│  │   Routes     │  │   Server     │  │     Auth     │  │
│  └──────────────┘  └──────────────┘  └──────────────┘  │
│                            │                            │
│                  ┌──────────────┐                       │
│                  │   SQLite DB  │                       │
│                  └──────────────┘                       │
└─────────────────────────────────────────────────────────┘
\end{verbatim}

\subsection{Database Schema}

The application uses five main tables:

\begin{enumerate}[leftmargin=*]
    \item \textbf{users:} User accounts with email, password hash, profile info
    \item \textbf{threads:} Discussion threads with categories and tags
    \item \textbf{replies:} Comments on threads (hierarchical structure)
    \item \textbf{likes:} User likes on threads and replies
    \item \textbf{notifications:} Activity notifications for users
    \item \textbf{messages:} Direct messages between users
    \item \textbf{message\_read\_status:} Read receipts for messages
\end{enumerate}

\section{Key Features Implementation}

\subsection{Custom JWT Authentication System}

Replaced third-party authentication (Clerk) with a custom implementation:

\begin{lstlisting}[language=JavaScript, caption=JWT Token Generation]
// backend/src/config/jwt.ts
export function generateToken(userId: number): string {
  return jwt.sign(
    { userId },
    process.env.JWT_SECRET || 'fallback-secret',
    { expiresIn: '7d' }
  );
}
\end{lstlisting}

\textbf{Key Components:}
\begin{itemize}[leftmargin=*]
    \item Password hashing with bcrypt (10 salt rounds)
    \item JWT token generation with 7-day expiration
    \item Custom middleware for route protection
    \item Frontend AuthProvider with React Context
    \item Automatic token injection in API requests
    \item Token storage in localStorage
\end{itemize}

\subsection{Database Migration Strategy}

Successfully migrated from PostgreSQL to SQLite:

\textbf{Challenges Solved:}
\begin{itemize}[leftmargin=*]
    \item Converted positional parameters (\$1, \$2) to ? placeholders
    \item Replaced PostgreSQL-specific functions:
    \begin{itemize}
        \item \texttt{NOW()} $\rightarrow$ \texttt{datetime('now')}
        \item \texttt{ILIKE} $\rightarrow$ \texttt{LIKE}
        \item \texttt{LEFT()} $\rightarrow$ \texttt{SUBSTR()}
    \end{itemize}
    \item Updated AUTO\_INCREMENT syntax for SQLite
    \item Modified query return values (insertId vs lastInsertRowid)
    \item Implemented triggers for timestamp updates
\end{itemize}

\subsection{Real-time Communication}

Implemented bidirectional communication using Socket.IO:

\begin{lstlisting}[language=JavaScript, caption=Socket.IO Integration]
// backend/src/realtime/io.ts
io.on('connection', (socket) => {
  console.log('[io connection]------>', socket.id);
  
  socket.on('disconnect', () => {
    console.log('[io disconnect]------>', socket.id);
  });
});
\end{lstlisting}

\textbf{Real-time Features:}
\begin{itemize}[leftmargin=*]
    \item Direct messaging with instant delivery
    \item Online presence indicators
    \item Live notification updates
    \item Message read status tracking
\end{itemize}

\subsection{File Upload System}

Replaced cloud storage (Cloudinary) with local file system:

\begin{itemize}[leftmargin=*]
    \item Multer disk storage configuration
    \item File type validation (images only)
    \item Automatic directory creation
    \item Static file serving via Express
    \item Image upload for threads and profiles
\end{itemize}

\section{API Design}

\subsection{RESTful Endpoints}

\subsubsection{Authentication Routes}
\begin{itemize}[leftmargin=*]
    \item \texttt{POST /api/auth/register} - User registration
    \item \texttt{POST /api/auth/login} - User login
    \item \texttt{GET /api/auth/me} - Get current user (protected)
\end{itemize}

\subsubsection{Thread Routes}
\begin{itemize}[leftmargin=*]
    \item \texttt{GET /api/threads} - List threads with filters
    \item \texttt{POST /api/threads} - Create new thread (protected)
    \item \texttt{GET /api/threads/:id} - Get thread details
    \item \texttt{POST /api/threads/:id/like} - Toggle like (protected)
    \item \texttt{POST /api/threads/:id/replies} - Add reply (protected)
\end{itemize}

\subsubsection{User Routes}
\begin{itemize}[leftmargin=*]
    \item \texttt{GET /api/users/:id} - Get user profile
    \item \texttt{PATCH /api/users/:id} - Update profile (protected)
\end{itemize}

\subsubsection{Chat Routes}
\begin{itemize}[leftmargin=*]
    \item \texttt{GET /api/chat/users} - List chat users
    \item \texttt{GET /api/chat/:userId} - Get conversation
    \item \texttt{POST /api/chat/:userId} - Send message (protected)
\end{itemize}

\subsubsection{Notification Routes}
\begin{itemize}[leftmargin=*]
    \item \texttt{GET /api/notifications} - List notifications (protected)
    \item \texttt{POST /api/notifications/:id/read} - Mark as read
    \item \texttt{GET /api/notifications/unread-count} - Get unread count
\end{itemize}

\subsection{Request/Response Format}

All API responses follow a consistent JSON structure:

\begin{lstlisting}[language=JSON, caption=Success Response]
{
  "id": 1,
  "email": "user@example.com",
  "displayName": "John Doe",
  "handle": "@john",
  "createdAt": "2026-02-20T12:00:00.000Z"
}
\end{lstlisting}

\begin{lstlisting}[language=JSON, caption=Error Response]
{
  "error": {
    "message": "Route not found",
    "status": 404
  }
}
\end{lstlisting}

\section{Security Implementation}

\subsection{Authentication Security}
\begin{itemize}[leftmargin=*]
    \item Passwords hashed with bcrypt (10 rounds)
    \item JWT tokens with expiration (7 days)
    \item Protected routes with authentication middleware
    \item Token validation on every protected request
    \item Automatic logout on 401 responses
\end{itemize}

\subsection{Input Validation}
\begin{itemize}[leftmargin=*]
    \item Email format validation
    \item Password strength requirements
    \item File type validation for uploads
    \item SQL injection prevention via parameterized queries
\end{itemize}

\subsection{Best Practices}
\begin{itemize}[leftmargin=*]
    \item Environment variables for sensitive data
    \item CORS configuration
    \item Error handling middleware
    \item Secure HTTP headers
\end{itemize}

\section{Frontend Architecture}

\subsection{Component Structure}
\begin{verbatim}
src/
├── app/                    # Next.js App Router pages
│   ├── page.tsx           # Home (threads list)
│   ├── sign-in/           # Login page
│   ├── sign-up/           # Registration page
│   ├── threads/
│   │   ├── new/           # Create thread
│   │   └── [id]/          # Thread details
│   ├── chat/              # Direct messages
│   ├── notifications/     # Notifications page
│   └── profile/           # User profile
├── components/            # Reusable components
│   ├── layout/
│   │   └── navbar.tsx     # Navigation bar
│   ├── threads/
│   │   └── threads-home.tsx
│   ├── chat/
│   │   └── direct-chat-panel.tsx
│   └── ui/                # shadcn/ui components
├── contexts/
│   └── auth-context.tsx   # Authentication state
├── hooks/
│   ├── use-socket.ts      # Socket.IO hook
│   └── use-notification-count.tsx
├── lib/
│   ├── api-client.ts      # Axios instance
│   └── utils.ts           # Utility functions
└── types/                 # TypeScript definitions
\end{verbatim}

\subsection{State Management}

Custom React Context for authentication:

\begin{lstlisting}[language=JavaScript, caption=Auth Context Provider]
export function AuthProvider({ children }: { children: ReactNode }) {
  const [user, setUser] = useState<User | null>(null);
  const [loading, setLoading] = useState(true);
  
  useEffect(() => {
    const token = localStorage.getItem('token');
    if (token) {
      fetchCurrentUser();
    } else {
      setLoading(false);
    }
  }, []);
  
  return (
    <AuthContext.Provider value={{ user, login, register, logout }}>
      {children}
    </AuthContext.Provider>
  );
}
\end{lstlisting}

\section{Technical Challenges \& Solutions}

\subsection{Challenge 1: Third-Party API Removal}

\textbf{Problem:} Application heavily relied on Clerk (auth) and Cloudinary (storage)

\textbf{Solution:}
\begin{itemize}[leftmargin=*]
    \item Implemented custom JWT authentication system
    \item Created local file storage with Multer
    \item Built custom React auth context
    \item Updated all components to use custom auth
\end{itemize}

\subsection{Challenge 2: Database Migration}

\textbf{Problem:} PostgreSQL syntax incompatible with SQLite

\textbf{Solution:}
\begin{itemize}[leftmargin=*]
    \item Analyzed all SQL queries across codebase
    \item Created systematic conversion strategy
    \item Updated 22+ queries across 4 repository files
    \item Implemented SQLite-specific datetime handling
\end{itemize}

\subsection{Challenge 3: Type Safety}

\textbf{Problem:} Maintaining type consistency across stack

\textbf{Solution:}
\begin{itemize}[leftmargin=*]
    \item Shared TypeScript interfaces for data models
    \item Custom type definitions for API responses
    \item Strict TypeScript configuration
    \item Runtime type validation where needed
\end{itemize}

\section{Performance Optimizations}

\begin{itemize}[leftmargin=*]
    \item \textbf{Database Indexing:} Unique constraints on email and handle
    \item \textbf{React Memoization:} useMemo for API client instances
    \item \textbf{Lazy Loading:} Component-level code splitting
    \item \textbf{Efficient Queries:} Optimized SQL with proper JOINs
    \item \textbf{WebSocket Efficiency:} Event-based real-time updates
\end{itemize}

\section{Testing \& Debugging}

\subsection{Testing Strategy}
\begin{itemize}[leftmargin=*]
    \item Manual testing of all user flows
    \item API testing via frontend integration
    \item Real-time feature testing with multiple clients
    \item Error scenario testing (network failures, auth errors)
\end{itemize}

\subsection{Debugging Approach}
\begin{itemize}[leftmargin=*]
    \item Console logging for development
    \item TypeScript for compile-time error catching
    \item Browser DevTools for frontend debugging
    \item Network tab for API inspection
\end{itemize}

\section{Deployment Considerations}

\subsection{Environment Variables}
\begin{lstlisting}[language=bash, caption=Backend .env]
PORT=5000
JWT_SECRET=your-secret-key
DATABASE_PATH=./database.sqlite
UPLOAD_DIR=./uploads
\end{lstlisting}

\subsection{Production Readiness}
\begin{itemize}[leftmargin=*]
    \item Environment-based configuration
    \item Error handling middleware
    \item Logging infrastructure ready
    \item Database migrations system
    \item Docker Compose configuration available
\end{itemize}

\section{Future Enhancements}

\subsection{Potential Improvements}
\begin{itemize}[leftmargin=*]
    \item \textbf{Admin Dashboard:} User management and moderation tools
    \item \textbf{Email Verification:} Confirm user email addresses
    \item \textbf{Password Reset:} Forgot password functionality
    \item \textbf{Rate Limiting:} Prevent API abuse
    \item \textbf{Caching:} Redis for session management
    \item \textbf{Search:} Full-text search for threads
    \item \textbf{Rich Text Editor:} Markdown or WYSIWYG support
    \item \textbf{File Attachments:} Support for non-image files
    \item \textbf{Thread Categories:} Better organization
    \item \textbf{User Mentions:} @ tagging in threads/replies
    \item \textbf{Reactions:} Multiple reaction types beyond likes
    \item \textbf{Analytics:} User engagement metrics
\end{itemize}

\section{Interview Discussion Points}

\subsection{Technical Skills Demonstrated}
\begin{itemize}[leftmargin=*]
    \item Full-stack TypeScript development
    \item RESTful API design and implementation
    \item Real-time bidirectional communication
    \item Database design and migration
    \item Authentication and authorization
    \item State management in React
    \item Modern frontend development (Next.js)
    \item Problem-solving and debugging
\end{itemize}

\subsection{Architectural Decisions}
\begin{enumerate}[leftmargin=*]
    \item \textbf{Why SQLite?} 
    \begin{itemize}
        \item Simplified deployment (no separate DB server)
        \item Sufficient for small to medium applications
        \item Zero configuration required
        \item Perfect for development and testing
    \end{itemize}
    
    \item \textbf{Why Custom Auth?}
    \begin{itemize}
        \item No external dependencies or costs
        \item Full control over authentication flow
        \item Better understanding of security
        \item Customizable to specific needs
    \end{itemize}
    
    \item \textbf{Why Next.js?}
    \begin{itemize}
        \item Server-side rendering capabilities
        \item Built-in routing system
        \item Excellent developer experience
        \item Production-ready optimizations
    \end{itemize}
    
    \item \textbf{Why Socket.IO?}
    \begin{itemize}
        \item Reliable WebSocket implementation
        \item Automatic fallback mechanisms
        \item Easy event-based API
        \item Great documentation and community
    \end{itemize}
\end{enumerate}

\subsection{Lessons Learned}
\begin{itemize}[leftmargin=*]
    \item Importance of database abstraction layers
    \item Value of TypeScript in preventing runtime errors
    \item Strategic approach to dependency management
    \item Systematic debugging methodology
    \item Balance between third-party libraries and custom code
\end{itemize}

\section{Code Quality Metrics}

\begin{itemize}[leftmargin=*]
    \item \textbf{TypeScript Coverage:} 100\% (all files in TypeScript)
    \item \textbf{Component Reusability:} High (shadcn/ui components)
    \item \textbf{Code Organization:} Modular structure with clear separation
    \item \textbf{Error Handling:} Comprehensive error middleware
    \item \textbf{Documentation:} Clear file structure and naming conventions
\end{itemize}

\section{Running the Application}

\subsection{Prerequisites}
\begin{itemize}[leftmargin=*]
    \item Node.js 18+ installed
    \item npm or yarn package manager
\end{itemize}

\subsection{Installation Steps}

\begin{lstlisting}[language=bash, caption=Backend Setup]
cd backend
npm install
npm run dev  # Runs on port 5000
\end{lstlisting}

\begin{lstlisting}[language=bash, caption=Frontend Setup]
cd frontend
npm install
npm run dev  # Runs on port 4000
\end{lstlisting}

\subsection{Access Points}
\begin{itemize}[leftmargin=*]
    \item Frontend: \url{http://localhost:4000}
    \item Backend API: \url{http://localhost:5000}
\end{itemize}

\section{Conclusion}

This project demonstrates comprehensive full-stack development capabilities, including:

\begin{itemize}[leftmargin=*]
    \item Modern web technology proficiency
    \item Problem-solving and debugging skills
    \item System architecture design
    \item Security best practices
    \item Real-time communication implementation
    \item Database management and migration
    \item API design and development
\end{itemize}

The application is production-ready with room for enhancement, showcasing both technical competence and scalability awareness.

\end{document}
